% Source: https://artofproblemsolving.com/community/c3h2779594_2016_mbmt_guts_round__p1_p15__montgomery_blair_math_tournament
% Stable identity: aops-topic-2779594-q12
The number $n=1$ qualifies directly. For $n>1$, pairing each divisor $d$ with $\dfrac{n}{d}$ gives divisor product $n^{\dfrac{\tau(n)}{2}}$, so $\tau(n)=12$. Possible exponent patterns are $(11),(5,1),(3,2),(2,1,1)$. Their smallest values are respectively $2^{11}=2048$, $2^5\cdot3=96$, $2^3\cdot3^2=72$, and $2^2\cdot3\cdot5=60$. In the last pattern, the only value below 72 is 60: keeping the squared prime 2, the next product of two other distinct primes is $3\cdot7$, giving 84; squaring 3 or a larger prime gives at least 90. Thus the three smallest overall values are 1,60,72, whose sum is $\boxed{133}$.
